\documentclass[11pt]{article}
\usepackage[a4paper,margin=2.5cm]{geometry}
\usepackage{amsmath,amssymb,amsthm,mathtools,booktabs}
\newtheorem{theorem}{Theorem}
\newtheorem{proposition}{Proposition}
\newtheorem{lemma}{Lemma}
\newtheorem{corollary}{Corollary}
\newtheorem{assumption}{Assumption}
\newtheorem{remark}{Remark}
\newcommand{\R}{\mathbb R}
\newcommand{\E}{\mathbb E}
\newcommand{\op}{\mathrm{op}}
\newcommand{\tr}{\operatorname{tr}}
\newcommand{\cond}{\operatorname{cond}}
\title{Encoder--Decoder Generalization with Heavy--Ball, Chebyshev, and PCG}
\author{}
\date{}

\begin{document}
\maketitle

\section{Setup and the correct error decomposition}

One raw PDE prompt is encoded into weak equations $(\widehat G,\widehat b)$.
The ideal weak system is $(G,b)$.  Put
\[
 H=\tau G^\top G+\lambda I,
 \quad c=\tau G^\top b,
 \quad z^\star=H^{-1}c,
\]
and define $\widehat H,\widehat c,\widehat z^\star$ analogously.  Here
$\tau,\lambda>0$.  A depth-$L$ controller returns $z_L$ for the encoded
system.  The first-principles decomposition is
\begin{equation}
 z_L-z^\star
 =\underbrace{z_L-\widehat z^\star}_{\text{finite-depth solver}}
 +\underbrace{\widehat z^\star-z^\star}_{\text{encoder/system perturbation}}.
 \label{eq:decomposition}
\end{equation}
The physical reconstruction error from the learned operator dictionary is a
third, downstream term.  It must not be conflated with either term in
\eqref{eq:decomposition}.

\begin{lemma}[Exact weak-system perturbation bound]
\label{lem:ridge-perturbation}
Let $\Delta G=\widehat G-G$ and $\Delta b=\widehat b-b$.  Then
\begin{align}
 \|\widehat H-H\|_{\op}
 &\leq \tau(\|G\|_{\op}+\|\widehat G\|_{\op})
 \|\Delta G\|_{\op},\label{eq:delta-H}\\
 \|\widehat c-c\|_2
 &\leq\tau\bigl(
 \|\widehat G\|_{\op}\|\Delta b\|_2
 +\|\Delta G\|_{\op}\|b\|_2\bigr),\label{eq:delta-c}
\end{align}
and
\begin{equation}
 \|\widehat z^\star-z^\star\|_2
 \leq \lambda^{-1}\|\widehat c-c\|_2
 +\lambda^{-2}\|\widehat H-H\|_{\op}\|c\|_2.
 \label{eq:ridge-transfer}
\end{equation}
\end{lemma}
\begin{proof}
Expand
$\widehat G^\top\widehat G-G^\top G
=\widehat G^\top\Delta G+\Delta G^\top G$ and
$\widehat G^\top\widehat b-G^\top b
=\widehat G^\top\Delta b+\Delta G^\top b$.
For the last display use
\[
 \widehat H^{-1}\widehat c-H^{-1}c
 =\widehat H^{-1}(\widehat c-c)
 -\widehat H^{-1}(\widehat H-H)H^{-1}c
\]
and $H,\widehat H\succeq\lambda I$.
\end{proof}

This lemma is the missing bridge from encoder generalization to any of the
three controllers.  It uses the actual weak-feature errors rather than an
unjustified independence approximation between $G$ and the encoder residual.

\section{One learned geometry, three optimization terms}

Let $B_\theta(\mathcal C)\succ0$ be computed once from the encoded prompt and
kept fixed through depth.  Suppose
\[
 \sigma(B_\theta^{1/2}\widehat H B_\theta^{1/2})
 \subseteq[\mu,M],
 \qquad \kappa=M/\mu.
\]
All bounds below are conditional on the encoded prompt and can subsequently be
averaged over tasks.

\subsection{Heavy--Ball}

For the optimal constant coefficients, define
\[
 q=\frac{\sqrt\kappa-1}{\sqrt\kappa+1},
 \qquad C_L(q)=1+L(1+q).
\]
With $z_{-1}=z_0=0$, the exact error polynomial obeys
\[
 p_{-1}^{\rm HB}=p_0^{\rm HB}=1,
 \quad
 p_{\ell+1}^{\rm HB}(x)
 =(1+\beta-\alpha x)p_\ell^{\rm HB}(x)
 -\beta p_{\ell-1}^{\rm HB}(x),
\]
and
\begin{equation}
 \|z_L^{\rm HB}-\widehat z^\star\|_{B_\theta^{-1}}
 \leq C_L(q)q^L
 \|\widehat z^\star\|_{B_\theta^{-1}}.
 \label{eq:hb-envelope}
\end{equation}
The polynomial is independent of the right-hand side.  Conditional Gaussian
trace formulas therefore remain exact.

For prompt-adaptive HB, let
$[\widehat\mu,\widehat M]$ be predicted from $\widehat H$ and the recovered
subspace but not from $\widehat b$, and substitute it into the analytic
minimax formulas for $(\alpha,\beta)$.  Define the actual stability event
\[
 \mathcal E_{\rm Jury}^{\rm HB}
 =\{0\leq\beta<1,\ \alpha M<2(1+\beta)\}.
\]
On this event the scalar recurrence is stable on the true spectrum.  Full
coverage $[\mu,M]\subseteq[\widehat\mu,\widehat M]$ is a stronger sufficient
condition giving equation~\eqref{eq:hb-envelope} with the predicted interval,
but it is not necessary for stability.  The resulting polynomial remains
independent of $\widehat b$, so the conditional trace calculation is
unchanged.  Across tasks, the spectral MLP and an additive failure term
proportional to $\Pr((\mathcal E_{\rm Jury}^{\rm HB})^c)$ enter the population
bound.

\subsection{Chebyshev}

Let $[\widehat\mu,\widehat M]$ be produced from $\widehat H$ and the recovered
subspace, but not from $\widehat b$.  On the coverage event
\[
 \mathcal E_{\rm cov}
 =\{[\mu,M]\subseteq[\widehat\mu,\widehat M]\},
\]
the fixed scalar recurrence generates a polynomial
$p_L^{\rm Ch}(x)$ independent of the right-hand side and gives
\begin{equation}
 \|z_L^{\rm Ch}-\widehat z^\star\|_{B_\theta^{-1}}
 \leq 2q_{\rm Ch}^{L}
 \|\widehat z^\star\|_{B_\theta^{-1}},
 \quad
 q_{\rm Ch}
 =\frac{\sqrt{\widehat M/\widehat\mu}-1}
 {\sqrt{\widehat M/\widehat\mu}+1}.
 \label{eq:cheb-envelope}
\end{equation}
Outside $\mathcal E_{\rm cov}$ there is no convergence guarantee.  A valid
population-risk statement must therefore include either a deterministic
certificate or a sensitivity-weighted failure term.  More precisely, if
$S(\mathcal C)$ is any taskwise envelope for the squared physical error when
coverage fails, then the uncontrolled contribution is
\begin{equation}
 \mathfrak F_{\rm cov}
 =\E\!\left[S(\mathcal C)\mathbf 1_{\mathcal E_{\rm cov}^{c}}\right],
 \label{eq:sensitivity-coverage-failure}
\end{equation}
not merely a constant times $\Pr(\mathcal E_{\rm cov}^{c})$ unless $S$ is
uniformly bounded.  This distinction is necessary because a prompt-independent
interval can attain high marginal coverage while missing a small set of
spectrally sensitive tasks.  Reporting only interval-regression MSE or marginal
coverage is therefore insufficient.

The scalar recurrence is also a strict tied loop.  If
$z_{-1}=z_0=0$, its state is
$[z_\ell,z_{\ell-1},\alpha_\ell,\beta_\ell]$, and the next block performs the
exact hard-coded update
\[
 z_{\ell+1}=z_\ell+\alpha_\ell
 (\widehat b-\widehat H z_\ell)
 +\beta_\ell(z_\ell-z_{\ell-1}).
\]
For fixed predicted endpoints, the scalar coefficient state has a closed-form
solution and may equivalently be compiled once into a length-$L$ schedule.
This compilation changes neither the polynomial nor the hypothesis class; it
does not turn the controller into a non-looping architecture.

To call the interval head an in-context controller rather than a global
calibration artefact, let $\xi(\mathcal C)$ be its prompt-derived features and
$h_\phi(\xi)=(\widehat\mu,\widehat M)$.  At test time $\phi$ is fixed and only
$\xi$ changes.  A necessary causal audit compares
\begin{equation}
 h_\phi(\xi(\mathcal C)),\qquad
 h_\phi(\xi(\pi(\mathcal C))),\qquad
 (\overline\mu,\overline M),
 \label{eq:icl-causal-audit}
\end{equation}
on identical held-out tasks, where $\pi$ is an independent task permutation
and $(\overline\mu,\overline M)$ is calibrated on disjoint tasks.  A material
loss of solver accuracy under permutation, beyond any change in aggregate
coverage, proves that the frozen head uses task-specific prompt information.
It does not by itself prove optimality, but it rules out the claim that the
gain is solely due to a learned global interval.

\subsection{PCG}

Fixed-preconditioner PCG satisfies
\begin{equation}
 \|z_L^{\rm PCG}-\widehat z^\star\|_{\widehat H}
 \leq 2q^L\|\widehat z^\star\|_{\widehat H}.
 \label{eq:pcg-envelope}
\end{equation}
It also minimizes the energy error over the depth-$L$ preconditioned Krylov
space.  Unlike the other two controllers, its realized polynomial depends on
$\widehat b$ through the adaptive inner products.  Thus the uniform envelope
transfers, but a fixed-polynomial Gaussian trace identity does not.

\begin{theorem}[Common encoder--decoder transfer]
\label{thm:common-transfer}
Let $\varepsilon_L(\mathcal C)$ denote the appropriate conditional solver
envelope from \eqref{eq:hb-envelope}, \eqref{eq:cheb-envelope}, or
\eqref{eq:pcg-envelope}, converted to Euclidean norm.  Then
\begin{align}
 \E\|z_L-z^\star\|_2^2
 \leq{}&2\E\bigl[\varepsilon_L(\mathcal C)^2
 \|\widehat z^\star\|_2^2\bigr]\notag\\
 &+4\lambda^{-2}\E\|\widehat c-c\|_2^2
 +4\lambda^{-4}\E\bigl[
 \|\widehat H-H\|_{\op}^2\|c\|_2^2\bigr].
 \label{eq:common-transfer}
\end{align}
For Chebyshev the first expectation is restricted to
$\mathcal E_{\rm cov}$, plus the sensitivity-weighted term
$\mathfrak F_{\rm cov}$ from
\eqref{eq:sensitivity-coverage-failure} after composition with the physical
decoder.
\end{theorem}
\begin{proof}
Apply $(a+b)^2\leq2a^2+2b^2$ to
\eqref{eq:decomposition}, insert Lemma~\ref{lem:ridge-perturbation}, and apply
the corresponding conditional solver bound.
\end{proof}

The same proof composes with a Lipschitz physical map
$z\mapsto A(z)\mapsto u_\star$ and with the operator-subspace projector error.
That last composition must be performed in a fixed orthonormal gauge or with a
covariant coefficient metric.

\section{Across-task generalization complexity}

Let $N$ be the number of independent training tasks.  The learned hypothesis
class contains the encoder and the common subspace/preconditioner head.
Controller-dependent additions are as follows.

\begin{center}
\begin{tabular}{p{0.13\linewidth}p{0.24\linewidth}p{0.28\linewidth}p{0.22\linewidth}}
\toprule
Controller & Additional learned class & Required stability event & Trace linear in $b$\\
\midrule
HB & shared scalars or spectral-interval MLP & Jury margin or interval coverage & yes, if bounds ignore $b$\\
Chebyshev & spectral-interval MLP & interval coverage with slack & yes, if bounds ignore $b$\\
PCG & none & positive denominators or converged mask & no\\
\bottomrule
\end{tabular}
\end{center}

For tied-scalar HB, compact stable parameters add only two coordinates to a
covering argument.  Adaptive HB and Chebyshev instead add the interval MLP's
actual norm-controlled parameter complexity.  For adaptive HB the Lipschitz
constant grows as the Jury margin shrinks; for Chebyshev it grows as the
coverage slack shrinks.  For PCG, the solver adds no trainable parameter, but a pathwise
Lipschitz proof requires denominators bounded away from zero until the fixed
convergence mask activates.  The numerical mask is therefore part of the
mathematical architecture, not an implementation footnote.

Under bounded prompts, bounded network norms, a uniform stability margin, and
a Lipschitz truncated loss, standard parameter-covering or Rademacher arguments
retain an $N^{-1/2}$ leading rate.  The constants and entropy terms differ and
must not be hidden inside the old one-parameter Richardson bound.

For the Chebyshev interval head, this across-task generalization statement is
about the frozen map $h_\phi:\xi(\mathcal C)\mapsto
(\widehat\mu,\widehat M)$, not about approximating multiplication, division,
or the matrix--vector recurrence with an MLP.  Those algebraic operations are
exact architectural primitives.  The learned class contains only the
low-dimensional task-to-interval policy.  Its risk must combine the covered
solver term and $\mathfrak F_{\rm cov}$; an $N^{-1/2}$ bound for marginal
coverage alone does not control the latter without a bound on task
sensitivity or a sensitivity-weighted calibration loss.

\section{Joint predictive theory for controller and hyperparameters}
\label{sec:joint-predictive}

The preceding envelopes certify a controller chosen in advance.  They do not
yet predict which controller, depth, or coefficients minimize the physical
risk.  We now give an exact finite-dimensional objective before taking any
replica limit.

For one encoded task, write
\[
 A_\theta=B_\theta^{1/2}\widehat H B_\theta^{1/2},
 \qquad d_\theta=B_\theta^{1/2}\widehat c,
 \qquad y^\star=A_\theta^{-1}d_\theta,
 \qquad \widehat z^\star=B_\theta^{1/2}y^\star.
\]
Let $p_{s,L}(\,\cdot\,;\eta)$ be the exact residual polynomial for controller
$s\in\{\mathrm{HB},\mathrm{Ch},\mathrm{R}\}$, depth $L$, and scalar
hyperparameters $\eta$.  Starting from zero gives
\[
 z_{s,L}-\widehat z^\star
 =-B_\theta^{1/2}p_{s,L}(A_\theta;\eta)y^\star.
\]
Let $F_{\mathcal C}(z)$ be the physical query decoder, put
$r_{\mathcal C}=F_{\mathcal C}(\widehat z^\star)-u_\star$, and denote its
Jacobian at $\widehat z^\star$ by $J_{\mathcal C}$.  If
$A_\theta=U\operatorname{diag}(\lambda_i)U^\top$, define
\[
 v=U^\top y^\star,
 \qquad
 C=U^\top B_\theta^{1/2}J_{\mathcal C}^\top J_{\mathcal C}
 B_\theta^{1/2}U.
\]
Two taskwise signed spectral measures contain all quantities needed by the
linearized physical risk:
\begin{align*}
 \Gamma_{\mathcal C}
 &=\sum_{i,j}v_i v_j C_{ij}\,
 \delta_{(\lambda_i,\lambda_j)},\\
 \zeta_{\mathcal C}
 &=\sum_i v_i
 \left\langle r_{\mathcal C},
 J_{\mathcal C}B_\theta^{1/2}Ue_i\right\rangle
 \delta_{\lambda_i}.
\end{align*}
Although $\Gamma_{\mathcal C}$ need not be a positive measure, its quadratic
form is nonnegative for every real polynomial.  No commutation between the
physical Jacobian and $A_\theta$ has been assumed.

\begin{theorem}[Exact joint polynomial risk]
\label{thm:joint-polynomial-risk}
If $F_{\mathcal C}$ is affine in $z$, the population physical risk of any
right-hand-side-independent polynomial controller is exactly
\begin{align}
 \mathcal R_{s,L}(\theta,\eta)
 ={}&\mathcal R_{\rm enc}(\theta)
 -2\int p_{s,L}(\lambda;\eta)\,
 d\overline\zeta_\theta(\lambda)\notag\\
 &+\iint p_{s,L}(\lambda;\eta)
 p_{s,L}(\lambda';\eta)\,
 d\overline\Gamma_\theta(\lambda,\lambda'),
 \label{eq:exact-joint-risk}
\end{align}
where $\mathcal R_{\rm enc}=\E\|r_{\mathcal C}\|^2$ and the bars denote
expectation over tasks.  If $F_{\mathcal C}$ has a locally Lipschitz Jacobian,
the same display is the risk of its first-order physical linearization and
the omitted remainder is
$O(\E[\|r_{\mathcal C}\|\|e_{s,L}\|^2]
+\E\|e_{s,L}\|^3+\E\|e_{s,L}\|^4)$ under bounded moments, where
$e_{s,L}=z_{s,L}-\widehat z^\star$.
\end{theorem}
\begin{proof}
For an affine decoder, insert
$F_{\mathcal C}(z_{s,L})-u_\star
=r_{\mathcal C}-J_{\mathcal C}B_\theta^{1/2}
p_{s,L}(A_\theta)y^\star$ and expand its squared norm in the eigenbasis of
$A_\theta$.  The nonlinear statement follows from Taylor's theorem and
Cauchy--Schwarz.
\end{proof}

The predictive, latency-aware choice is consequently
\begin{equation}
 (s^\star,L^\star,\eta^\star,\theta^\star)
 \in\operatorname*{argmin}_{s,L,\eta,\theta}
 \left\{\mathcal R_{s,L}(\theta,\eta)
 +\tau_{\rm hw}T_s(L,\theta)\right\},
 \label{eq:joint-hyperparameter-predictor}
\end{equation}
subject to the controller's deterministic stability constraints.  Here
$T_s$ is a measured or analytic hardware cost and $\tau_{\rm hw}$ selects a
point on the risk--latency Pareto frontier.  Equation
\eqref{eq:joint-hyperparameter-predictor} predicts the preconditioner through
$\theta$, the controller, its depth, and its scalar hyperparameters jointly;
it is not a post-hoc convergence bound.

For HB, $\eta=(\alpha,\beta)$ and the feasible set is
\[
 0\leq\beta<1,
 \qquad
 \alpha\lambda_+(\theta)<2(1+\beta).
\]
Given the two averaged measures, the distribution-aware optimum is only a
two-dimensional deterministic minimization.  Its gradients require no
learned arithmetic: differentiate the exact recurrence according to
\begin{align*}
 \partial_\alpha p_{\ell+1}
 &=-\lambda p_\ell+(1+\beta-\alpha\lambda)
 \partial_\alpha p_\ell-\beta\partial_\alpha p_{\ell-1},\\
 \partial_\beta p_{\ell+1}
 &=p_\ell-p_{\ell-1}+(1+\beta-\alpha\lambda)
 \partial_\beta p_\ell-\beta\partial_\beta p_{\ell-1}.
\end{align*}
If only a support certificate $[\mu,M]$ is known, the distribution-free
minimax prediction is the closed form
\[
 \alpha_{\rm mm}=\frac{4}{(\sqrt M+\sqrt\mu)^2},
 \qquad
 \beta_{\rm mm}=
 \left(\frac{\sqrt M-\sqrt\mu}{\sqrt M+\sqrt\mu}\right)^2.
\]
Thus the closed-form theory and the distribution-aware theory are two levels
of the same predictor, not competing heuristics.

For a prompt-conditioned Chebyshev controller, the Bayes-optimal interval
policy is likewise defined without reference to a neural network:
\begin{equation}
 h^\star(\xi)
 \in\operatorname*{argmin}_{\widehat\mu,\widehat M}
 \E\!\left[\mathcal R_{\mathrm{Ch},L}
 (\widehat\mu,\widehat M)\mid\xi\right],
 \label{eq:bayes-chebyshev-policy}
\end{equation}
including the sensitivity-weighted failure cost from
\eqref{eq:sensitivity-coverage-failure}.  The small MLP is merely a
norm-controlled function class used to estimate $h^\star$ from finitely many
tasks.  Training it on endpoint MSE alone targets a different and generally
suboptimal object.

The assumptions used so far are only: SPD encoded normal operators, a
preconditioner and scalar policy fixed before seeing the query right-hand
side, differentiability of the physical decoder for the local nonlinear
statement, and the displayed moments.  Gaussianity, simultaneous
diagonalization, isotropy, and a large-system limit are not required.  If the
physical metric is the solver energy itself, $C$ is diagonal in the spectral
basis and $\Gamma$ reduces to a positive one-dimensional measure.

\subsection{What joint training can and cannot change}

Let $\mathcal L_L(\vartheta)$ be the training loss through a depth-$L$
polynomial decoder and $\mathcal L_\infty(\vartheta)$ the same loss through
the exact encoded solve, where $\vartheta$ denotes encoder and preconditioner
parameters.  On a compact uniformly stable parameter set, define
\[
 \delta_L=\sup_{\vartheta}\left(
 \|p_L(A_\vartheta)\|_{\op}
 +\|D_Ap_L(A_\vartheta)\|_{\op\to\op}\right).
\]

\begin{proposition}[Gradient fidelity of a deep fixed polynomial]
\label{prop:gradient-fidelity}
If the weak encoder, preconditioner, and physical loss have uniformly bounded
first derivatives on that set, then for a constant $K$ independent of $L$,
\[
 \|\nabla_\vartheta\mathcal L_L
 -\nabla_\vartheta\mathcal L_\infty\|
 \leq K\delta_L.
\]
For stable minimax HB or Chebyshev, $\delta_L$ is a polynomial factor in $L$
and $1/\mu$ times $q^L$; for optimal Richardson its exponential factor is
$(\kappa-1)/(\kappa+1)$ instead of
$(\sqrt\kappa-1)/(\sqrt\kappa+1)$.
\end{proposition}
\begin{proof}
Differentiate the exact identity
$z_L-\widehat z^\star=-B_\theta^{1/2}p_L(A_\theta)y^\star$ and apply the chain
rule.  Uniform derivative bounds control every factor except the polynomial
and its Fr\'echet derivative, which are precisely the two terms in
$\delta_L$.  The rates follow from the corresponding stable polynomial
recurrences and their differentiated recurrences.
\end{proof}

This proposition predicts that sufficiently deep HB and Richardson can learn
the same identifiable operator subspace even though HB reaches a prescribed
gradient fidelity at a smaller asymptotic depth.  The controller advantage
appears when depth or latency is binding, or on spectral tails.  It should not
be misreported as a different identifiability theorem for the encoder.

\subsection{An equivariant one-head closure with no eigenvector order parameter}

The coordinate-token Ritz head used in the first experiments is not invariant
under an orthogonal change of coefficient dictionary.  Its replica closure
therefore needs coordinate/eigenvector overlaps in addition to eigenvalues.
The following alternative removes that hidden assumption.  All Ritz algebra
is exact; one softmax head learns only how to allocate correction mass among
the spectral tokens.

Let $\widehat H=V\operatorname{diag}(h_1,\ldots,h_K)V^\top\succ0$, fix a certified
$\overline L>1$, and set
\[
 s_H=h_{\max}/\overline L,
 \qquad \lambda_i=h_i/s_H\in(0,\overline L].
\]
Let $x_i=x(\lambda_i;\nu_K)$ contain only permutation-invariant spectral
features, where $\nu_K=K^{-1}\sum_i\delta_{\lambda_i}$.  For slot $r$, define
\begin{equation}
\begin{aligned}
 \omega_{ri}
 &=\frac{\exp\{q_r^\top k_\theta(x_i)/\sqrt{d_h}\}}
 {\sum_j\exp\{q_r^\top k_\theta(x_j)/\sqrt{d_h}\}},\\
 g_i&=1-\prod_{r=1}^{S}(1-\omega_{ri}),\\
 B_\theta(\widehat H)
 &=V\operatorname{diag}\!\left(
 \frac{1-g_i}{s_H}+\frac{g_i}{h_i}
 \right)V^\top .
\end{aligned}
 \label{eq:equivariant-ritz-head}
\end{equation}
There is no MLP and no learned multiplication, reciprocal, eigensolver, or
recurrence in this definition.

\begin{theorem}[Exact finite-$K$ spectral closure]
\label{thm:equivariant-ritz-closure}
The map in \eqref{eq:equivariant-ritz-head} is SPD and orthogonally covariant:
for every orthogonal $Q$,
\[
 B_\theta(Q^\top\widehat H Q)=Q^\top B_\theta(\widehat H)Q.
\]
It is independent of the eigenbasis selected inside repeated eigenspaces.
Moreover,
\begin{equation}
\begin{aligned}
 A_\theta
 &=B_\theta^{1/2}\widehat H B_\theta^{1/2}
 =V\operatorname{diag}(a_1,\ldots,a_K)V^\top,\\
 a_i&=(1-g_i)\lambda_i+g_i,
 \qquad 0<a_i\leq\overline L,
 \qquad \sum_i g_i\leq S.
\end{aligned}
 \label{eq:effective-pushforward-finite}
\end{equation}
Hence the effective spectral law is known exactly from the raw eigenvalue law
and the one-head gates; no eigenvector replica order parameter is required.
In particular the HB Jury constraint holds for every prompt whenever
$\alpha\overline L<2(1+\beta)$.
\end{theorem}
\begin{proof}
Equal eigenvalues have equal invariant features and therefore equal gates, so
the spectral functional is well defined on the corresponding eigenspace and
is orthogonally covariant.  Multiplication of the commuting eigenvalue factors
gives $a_i$.  Since $g_i\in[0,1]$, $a_i$ lies between $\lambda_i$ and $1$.
Finally, the union bound gives
$g_i\leq\sum_r\omega_{ri}$ and summing over $i$ gives
$\sum_i g_i\leq S$.
\end{proof}

\begin{corollary}[exact-spectrum Chebyshev needs no interval MLP]
\label{cor:exact-head-chebyshev}
Under Theorem~\ref{thm:equivariant-ritz-closure}, the endpoints
\[
    \mu_{\mathcal C}=\min_i a_i,
    \qquad
    L_{\mathcal C}=\max_i a_i
\]
are already produced by the exact Ritz algebra of the preconditioner head.
Feeding these two scalars to the hard-coded Chebyshev coefficient recurrence
gives the exact promptwise Chebyshev controller that is minimax on the
enclosing interval $[\mu_{\mathcal C},L_{\mathcal C}]$, without an interval
MLP and without an additional eigendecomposition.  This does not claim
minimaxity on the finite discrete spectrum itself.  A learned interval head is
needed only when the spectral decomposition is not materialized, or when its
cost is deliberately amortized at larger coefficient dimension.
\end{corollary}

\begin{proposition}[a global covariance loses all directions under RRS]
\label{prop:rrs-global-covariance}
Let $H\succ0$ be any random prompt normal with
$\mathbb E\operatorname{tr}H<\infty$, and let $Q$ be an independent Haar
orthogonal matrix.  Then
\[
 \mathbb E_{H,Q}[Q^\top H Q]
 =\frac{\mathbb E\operatorname{tr}H}{K}I.
\]
Consequently the inverse-population-covariance preconditioner is a scalar
multiple of the identity and cannot retain a preferred eigenspace.  In
contrast, every promptwise orthogonally covariant head satisfies
\[
 B_\theta(Q^\top H Q)=Q^\top B_\theta(H)Q,
\]
so $B_\theta(Q^\top H Q)^{1/2}(Q^\top H Q)
B_\theta(Q^\top H Q)^{1/2}$ has the same eigenvalues as
$B_\theta(H)^{1/2}HB_\theta(H)^{1/2}$.  Thus its directional correction
survives arbitrary contextwise rotations without storing their orientation
in the weights.
\end{proposition}
\begin{proof}
Haar invariance makes $\mathbb E_Q[Q^\top H Q]$ commute with every orthogonal
matrix, hence it is scalar; trace preservation fixes the scalar.  The second
claim is Theorem~\ref{thm:equivariant-ritz-closure}, and orthogonal similarity
preserves the effective spectrum.
\end{proof}

This statement needs neither Gaussianity nor a deterministic limiting
spectrum.  It also exposes a stability distinction hidden by fixed-spectrum
models.  If the largest eigenvalue of $H$ has unbounded support, then for any
fixed scalar preconditioner, $\alpha>0$, and $\beta<1$, the HB Jury inequality
fails with positive probability.  No nonzero globally fixed step can be
uniformly certified without a spectral-tail assumption.  The promptwise
normalization in \eqref{eq:equivariant-ritz-head}, however, gives
$a_i\leq\overline L$ deterministically.  A finite-law risk or CVaR predictor
can still select the best global coefficients, but it must report this tail
risk rather than silently assuming bounded Gaussian quadratic forms.

This finite statement already closes the exact risk in
Theorem~\ref{thm:joint-polynomial-risk}.  A particularly simple scalable
limit is also available.  Suppose $\nu_K\Rightarrow\nu$ on a compact subset
of $(0,\overline L]$, $S/K\to\sigma$, the empirical slot-query law converges
to $\pi$, and the score
$\ell(q,\lambda;\nu)=q^\top k_\theta(x(\lambda;\nu))/\sqrt{d_h}$ is bounded
and continuous.  Then
\begin{equation}
\begin{aligned}
 Z(q)&=\int e^{\ell(q,t;\nu)}\,\nu(dt),\\
 g_{\theta,\sigma}(\lambda;\nu)
 &=1-\exp\!\left[-\sigma\int
 \frac{e^{\ell(q,\lambda;\nu)}}{Z(q)}\,\pi(dq)\right],\\
 \nu_{\rm eff}
 &=\bigl[(1-g_{\theta,\sigma})\lambda
       +g_{\theta,\sigma}\bigr]_{\#}\nu .
\end{aligned}
 \label{eq:softmax-head-limit}
\end{equation}
This follows from the softmax law of large numbers and
$\log(1-\omega)=-\omega+o(K^{-1})$.  It requires neither Gaussian prompts,
isotropic eigenvectors, nor commutation with the physical Jacobian.  If $S$
is fixed, then $\sigma=0$ and a bounded-score head cannot change the bulk
law: it is a finite spectral/outlier correction.  A nontrivial bulk
preconditioner requires $S=\Theta(K)$ or a concentrating low-temperature
attention regime.  This scaling distinction is part of the prediction, not
an experimental tuning rule.

Combining \eqref{eq:softmax-head-limit} with
\eqref{eq:joint-hyperparameter-predictor} predicts
$(\theta,\sigma,\alpha,\beta,L)$ jointly.  With only support information it
uses the closed-form minimax HB pair; with the limiting weighted law it uses
the two-scalar distribution-aware risk.  Thus the head and HB coefficients
are not calibrated in two unrelated stages.

\section{Replica and random-matrix substitutions}

Let $\nu_\theta$ be the limiting joint spectral law of the encoded normal
operator and the fixed learned preconditioner.  The preconditioner replica
order parameters must include at least the physical-space overlap and the
observable-space overlap
\[
 m_T=\frac1r\tr(P_{\widehat T}P_T),
 \qquad
 m_S=\frac1s\tr(P_{Q_\theta}P_{S(G)}),
\]
together with the projected Ritz spectral law.  A full learned matrix
$W_{\rm dec}$ and a recovered projector $Q_\theta Q_\theta^\top$ are not the
same order parameter.

For a commuting model with isotropic query covariance, Richardson's filter is
replaced by
\begin{align*}
 \mathcal R_{\rm solve,L}^{\rm HB}
 &\longrightarrow
 \int \frac{p_L^{\rm HB}(p\widehat a)^2}{\widehat a^2}
 \,d\nu_\theta(\widehat a,p),\\
 \mathcal R_{\rm solve,L}^{\rm Ch}
 &\longrightarrow
 \int \frac{p_L^{\rm Ch}(p\widehat a)^2}{\widehat a^2}
 \,d\nu_\theta(\widehat a,p),
\end{align*}
up to the query-variance scale.  These substitutions are valid because the
coefficients are conditionally independent of the right-hand side.

For adaptive HB or Chebyshev, let $\xi(\widehat H,B_\theta)$ denote the finite
vector of fixed spectral summaries supplied to the interval MLP.  If $\xi$
self-averages under the replica saddle, then the predicted endpoints and hence
the HB pair or Chebyshev schedule converge to deterministic functions of the
existing spectral order parameters.  The corresponding integral above is then
used with that limiting polynomial.  Without self-averaging, the joint law
must be enlarged to $\nu_\theta(\widehat a,p,\xi)$ and the appropriate filter
is averaged conditionally on $\xi$:
\[
 \mathcal R_{\rm solve,L}^{\rm Ch}
 \longrightarrow
 \int \frac{p_L^{\rm Ch}(p\widehat a;h_\phi(\xi))^2}
 {\widehat a^2}
 \,d\nu_\theta(\widehat a,p,\xi).
\]
The same joint law can evaluate the sensitivity-weighted coverage failure in
\eqref{eq:sensitivity-coverage-failure}.  No layer-indexed Krylov order
parameters are needed: once $\xi$ is fixed, both HB and Chebyshev are fixed
polynomials independent of the right-hand side.

For the genuinely joint prediction, the replica saddle must also determine
the averaged bivariate measure and cross measure from
Theorem~\ref{thm:joint-polynomial-risk}.  Let $q$ collect the encoder overlap,
preconditioner overlap, spectral law, and physical sensitivity order
parameters.  The required closure is
\[
 (\overline\Gamma_\theta,\overline\zeta_\theta,
 \mathcal R_{\rm enc}(\theta),T_s)
 \longrightarrow
 (\Gamma_q,\zeta_q,\mathcal R_{\rm enc}(q),T_s(q)).
\]
If hyperparameters affect representation learning, one must not optimize them
after solving a controller-independent encoder saddle.  Instead, with
$\Psi_{\rm tr}$ the replica training action, the predictive construction is
\begin{align}
 q_{s,L,\eta}
 &\in\operatorname*{extr}_q
 \Psi_{\rm tr}(q;s,L,\eta),
 \label{eq:controller-dependent-replica-saddle}\\
 (s^\star,L^\star,\eta^\star)
 &\in\operatorname*{argmin}_{s,L,\eta}
 \left\{\mathcal R_{s,L}^{\rm test}
 (q_{s,L,\eta},\eta)+\tau_{\rm hw}T_s(L,q_{s,L,\eta})\right\}.
 \label{eq:replica-hyperparameter-prediction}
\end{align}
The test risk in \eqref{eq:replica-hyperparameter-prediction} is obtained by
replacing the finite-task measures in \eqref{eq:exact-joint-risk} with
$(\Gamma_q,\zeta_q)$.  Equations
\eqref{eq:controller-dependent-replica-saddle}--
\eqref{eq:replica-hyperparameter-prediction} are the promised joint
encoder--decoder hyperparameter theory.  A frozen-encoder two-scalar HB fit is
the special case in which $q$ is fixed; the full theory predicts when that
modular approximation is valid through Proposition~\ref{prop:gradient-fidelity}.

There is no analogous fixed-filter substitution for PCG.  A PCG replica or
state evolution needs layer-indexed overlaps, for example residual norm,
preconditioned residual norm, direction norm, and their correlations with the
teacher, together with the scalar ratios $\alpha_\ell,\beta_\ell$.  Treating
PCG as Richardson with a better deterministic polynomial is incorrect.

\section{Optimizer boundary: AdamW versus Muon}

The collapsed gradient-flow ODEs previously derived for a factorized linear
attention encoder describe Euclidean gradient flow (and, approximately, small
step SGD/Adam regimes).  Muon replaces a matrix momentum update by an
approximately orthogonalized Newton--Schulz update.  Hence it changes the
training trajectory and implicit bias; the old ODE and its plateau times
cannot be reused.

The hypothesis class and a uniform generalization bound are optimizer-agnostic
once a parameter set is fixed, but the replica-selected saddle and dynamical
order parameters can depend on the optimizer.  Consequently Muon remains an
empirical ablation on large hidden encoder matrices.  It is not included in
the present replica claim, and AdamW/gradient flow remains the theoretical
reference.

\section{Audit of the Richardson-era claims}

\begin{center}
\begin{tabular}{p{0.34\linewidth}p{0.18\linewidth}p{0.40\linewidth}}
\toprule
Old ingredient & Status & Required replacement\\
\midrule
Encoder weak-feature replica law & retained & Feed its $\Delta G,\Delta b$ risk into
Lemma~\ref{lem:ridge-perturbation}.\\
One-step linear map $W_{\rm dec}s$ & insufficient & Separate recovered subspace,
factorized $B_\theta$, and recurrent controller.\\
Richardson spectral filter & retained structurally & Substitute the exact HB or
Chebyshev polynomial only.\\
Same substitution for PCG & invalid & Introduce adaptive Krylov trajectory order
parameters.\\
Linear-attention gradient-flow ODE under Muon & invalid & New optimizer-specific
dynamics or empirical-only claim.\\
Coefficient error without gauge fixing & invalid & Orthonormal operator dictionary
or covariant metric.\\
\bottomrule
\end{tabular}
\end{center}

This audit leaves the encoder replica calculation useful, but it prevents it
from being cited as a completed theory of the three recurrent decoders.

\end{document}
